% Revue de littérature : la taxe sur la valeur foncière
% Compilation : pdflatex lvt-revue.tex (deux passes) depuis fiscalite/lvt/
\documentclass[11pt,a4paper]{article}
% Préambule commun au brief et à la revue de littérature sur la taxe sur la valeur foncière.
% Même composition que le cours (cours/refaire-l-impot.tex) ; figures « nues » produites par
%   FIG_NU=1 python lvt/graphiques_lvt.py
\usepackage[T1]{fontenc}
\usepackage[utf8]{inputenc}
\usepackage{lmodern}
\usepackage[french]{babel}
\usepackage[a4paper,left=2.8cm,right=2.8cm,top=2.6cm,bottom=2.8cm]{geometry}
\usepackage{graphicx}
\usepackage[export]{adjustbox}
\usepackage{booktabs}
\usepackage{tabularx}
\usepackage{array}
\usepackage{xcolor}
\usepackage{titlesec}
\usepackage[labelfont=bf,labelsep=period,font=small,singlelinecheck=false,
            justification=raggedright,skip=6pt]{caption}
\usepackage{fancyhdr}
\usepackage{enumitem}
\usepackage{microtype}
\usepackage{textcomp}
\usepackage{amsmath}
\providecommand{\euro}{\texteuro}
\usepackage[hidelinks]{hyperref}

\definecolor{encre}{HTML}{1A1A19}
\definecolor{bleu}{HTML}{003399}
\definecolor{vert}{HTML}{1F7A4D}
\definecolor{ocre}{HTML}{C79100}
\definecolor{brique}{HTML}{993333}
\definecolor{filet}{HTML}{C8C6BE}

\titleformat{\section}{\normalfont\large\bfseries\color{bleu}}{\thesection.}{0.6em}{}
\titleformat{\subsection}{\normalfont\normalsize\bfseries\color{bleu}}{\thesubsection}{0.6em}{}
\titlespacing*{\section}{0pt}{2.4ex plus 1ex minus .2ex}{1.2ex plus .2ex}
\titlespacing*{\subsection}{0pt}{2ex plus 1ex minus .2ex}{0.8ex plus .2ex}

\pagestyle{fancy}
\fancyhf{}
\renewcommand{\headrulewidth}{0.4pt}
\fancyhead[R]{\small\thepage}

\setlength{\parindent}{0pt}
\setlength{\parskip}{0.62em}
\linespread{1.05}
\renewcommand{\arraystretch}{1.2}
\setlist{itemsep=2pt,topsep=2pt,parsep=0pt,leftmargin=1.4em}

\newenvironment{encadre}[1]
  {\par\medskip\noindent\begin{minipage}{\linewidth}%
   \color{encre}\setlength{\parskip}{0.5em}%
   \hrule height 0.8pt width \linewidth \vspace{6pt}%
   {\footnotesize\bfseries\color{vert}\MakeUppercase{#1}}\par\vspace{-2pt}%
   \small}
  {\par\vspace{4pt}\hrule height 0.4pt width \linewidth\end{minipage}\par\medskip}

\newcommand{\figbloc}[3]{%
  \begin{figure}[htbp]\centering
  \caption{#2}\label{#3}
  \includegraphics[max width=\linewidth,max height=0.40\textheight]{../figures/nu/#1.pdf}
  \IfFileExists{../figures/nu/#1_note.tex}{\par\vspace{4pt}%
    \parbox{\linewidth}{\raggedright\scriptsize\color{encre}\setlength{\parskip}{2pt}%
    \input{../figures/nu/#1_note.tex}}}{}%
  \end{figure}}

\newcommand{\pc}{\,\%}

\fancyhead[L]{\small\itshape Taxe sur la valeur foncière : revue de littérature}
\newcommand{\lu}[1]{{\footnotesize\color{ocre}[#1]}}

\begin{document}
\renewcommand{\figurename}{Figure}\renewcommand{\tablename}{Tableau}

\thispagestyle{empty}
\vspace*{0.4cm}
{\footnotesize\color{ocre}\bfseries REVUE DE LITTÉRATURE\quad·\quad OCTOBRE 2026\par}
\vspace{8pt}
{\fontsize{21}{25}\selectfont\bfseries\color{bleu}La taxe sur la valeur foncière : ce que l'on sait\par}
\vspace{8pt}
{\small Colin Baget. Pour chaque étude : la question, les données, ce qui identifie le résultat,
le résultat chiffré et ses limites. La mention entre crochets indique ce qui a été lu : texte
intégral, version de travail, ou résumé seulement. Les références dont l'existence ou le contenu
n'a pas pu être vérifié ne sont pas citées.\par}
\vspace{6pt}
{\color{bleu}\hrule height 1.2pt}

\section*{Où en est la littérature}

La théorie est ancienne et largement favorable ; les preuves empiriques sont rares, souvent
fragiles, et les plus récentes nuancent la théorie sur le point qui fonde son attrait. Quatre
résultats résument l'état des connaissances.

Le premier est théorique et solide : un impôt sur le seul terrain ne modifie aucune décision
de production si le terrain est évalué indépendamment de son usage, et son coût tombe sur les
propriétaires au moment de l'annonce par une baisse du prix du sol. Deux réserves théoriques
existent depuis les années 1970, sur l'incidence quand l'épargne peut se reporter du sol vers le
capital (Feldstein) et sur la neutralité quand l'évaluation dépend de l'usage (Bentick, Mills).

Le deuxième porte sur l'incidence et il est disputé. Les impôts fonciers ordinaires sont
largement capitalisés dans les prix (Koster et Pinchbeck, 2022) et répercutés sur les loyers
quand le bâti peut s'ajuster (Löffler et Siegloch, 2024). Pour une taxe sur le seul sol, les deux
études de la même réforme danoise se contredisent : capitalisation complète pour Høj et al. (2018),
aucune pour Nielsson et al. (2024), qui expliquent l'écart.

Le troisième porte sur la construction : les villes de Pennsylvanie qui taxent plus le terrain
que le bâti construisent un peu plus et plus dense, avec des effets modestes et souvent
statistiquement fragiles.

Le quatrième, le mieux établi, ne concerne pas la LVT mais ce qu'elle pourrait remplacer : les
droits de mutation réduisent fortement les transactions et la mobilité, avec des pertes sèches
élevées par euro prélevé.

\section{Théorie}

\textbf{Henry George et la rente foncière.} Arnott, R. et J. Stiglitz (1979), « Aggregate Land
Rents, Expenditure on Public Goods, and Optimal City Size », \emph{QJE} 93(4), 471-500
\lu{texte intégral}. Dans une ville de taille optimale dotée d'un bien public local pur, la somme
des rentes foncières différentielles est égale à la dépense en bien public : un impôt unique sur
ces rentes la finance exactement. Le résultat exige une grande économie, une allocation
Pareto-optimale et des rentes bien définies ; si l'une des conditions manque, l'impôt rapporte
trop ou trop peu. Avec congestion, rentes et recettes d'un péage optimal financent ensemble le bien
public. Les auteurs doutent que la concurrence produise des villes de taille optimale, donc que le
théorème décrive les villes observées.

\textbf{Incidence quand terre et capital sont substituables.} Feldstein, M. (1977), « The
Surprising Incidence of a Tax on Pure Rent: A New Answer to an Old Question », \emph{JPE} 85(2),
349-360 \lu{résumé et sources secondaires}. Dans un modèle à générations imbriquées où les ménages
épargnent pour leur retraite en détenant terre ou capital, la taxe sur la rente foncière réduit la
valeur de la terre ; l'épargne se reporte sur le capital, le taux d'intérêt baisse, la rente brute
de la terre et les salaires montent. La taxe est « au moins partiellement » transférée et le prix
de la terre peut même augmenter. Calvo, Kotlikoff et Rodriguez (1979, \emph{JPE} 87(4)) montrent
que le résultat disparaît avec un motif de legs altruiste.

\textbf{Neutralité vis-à-vis du calendrier de construction.} Bentick (1979, \emph{JPE} 87(4)) et
Mills, D. (1981), « The Non-Neutrality of Land Value Taxation », \emph{National Tax Journal} 34(1),
125-129 \lu{Mills : texte intégral} : une taxe sur la valeur capitalisée du terrain, si la valeur
est estimée à l'usage courant, frappe les revenus futurs avant qu'ils soient perçus et équivaut à
un relèvement du taux d'actualisation ; elle favorise les usages qui rapportent tôt. Selon Oates et Schwab, Tideman
(1982) et Wildasin (1982, \emph{NTJ} 35(1)) répondent que la neutralité est rétablie si le terrain
est évalué à son meilleur usage, indépendamment de ce qu'en fait le propriétaire. Oates et Schwab
(1997) donnent l'exemple numérique le plus clair et rappellent que l'idée selon laquelle la taxe
« pénalise » les terrains vides est fausse : elle est neutre parce qu'elle pèse autant que le
terrain soit bâti ou non.

\textbf{Étalement urbain.} Brueckner, J. et H.-A. Kim (2003), « Urban Sprawl and the Property
Tax », \emph{International Tax and Public Finance} 10(1), 5-23 \lu{résumé}. L'impôt sur la
propriété réduit le capital par hectare, ce qui étale la ville, mais aussi la taille des
logements, ce qui la densifie. Le second effet l'emporte quand l'élasticité de substitution entre
logement et autres biens est élevée ; l'impôt n'étale la ville que dans le cas, peu réaliste, de
préférences sans substitution. Remplacer l'impôt sur la propriété par une LVT ne densifie donc pas
mécaniquement : la question est empirique.

\textbf{Le retour du foncier et l'impôt optimal.} Bonnet, O., G. Chapelle, A. Trannoy et
É. Wasmer (2021), « Land is back, it should be taxed, it can be taxed », \emph{European Economic
Review} 134, 103696 \lu{texte intégral, post-print}. Deux contributions. Empirique : la hausse du
rapport patrimoine/revenu en France et au Royaume-Uni vient du logement et surtout de sa
composante foncière ; la part du foncier dans la valeur du parc français passe d'environ
16-20~\% des années 1980 à la fin des années 1990 à environ 50~\% vers 2007-2011 (lecture
graphique). Théorique : dans un modèle de Judd-Chamley avec terrain, structures et marché
locatif, un impôt foncier positif sans impôt sur le capital atteint le premier rang, en état
stationnaire comme en transition ; sans impôt foncier, une taxe sur les loyers combinée à une
subvention à l'investissement locatif fait presque aussi bien. Le travail a reçu le prix Maurice
Allais 2023. Il ne chiffre pas de recette pour la France.

Stiglitz, J. (2015), « New Theoretical Perspectives on the Distribution of Income and Wealth among
Individuals: Part IV. Land and Credit », NBER WP 21192 \lu{texte intégral}. Une part importante de
la hausse du rapport patrimoine/revenu vient de la valeur du sol, qui n'est pas du capital
productif. Une taxe sur le rendement du sol, plus-values comprises, réoriente l'épargne vers le
capital, élève le rapport capital/travail et les salaires, et réduit les inégalités.

Kumhof, M., N. Tideman, M. Hudson et C. Goodhart (2021), « Post-Corona Balanced-Budget
Super-Stimulus: The Case for Shifting Taxes onto Land », CEPR DP 16652 \lu{introduction et
résultats}. Modèle de croissance calibré sur les États-Unis avec capital, travail et sol. Porter
la taxe foncière de 0,55~\% à 5,55~\% de la valeur du sol sur vingt ans, en baissant de 28 points
l'impôt sur les revenus du capital et de 10 points celui sur le travail, accroîtrait le PIB de
14,8~\% et le bien-être de 3,4~\% de la consommation ; un impôt sur la fortune de même rendement
ferait environ moitié moins. Aucune validation empirique ; les auteurs qualifient eux-mêmes leurs
chiffres d'« extrêmement élevés » au regard de la littérature. Pas de publication en revue
trouvée.

\section{Capitalisation et incidence}

\textbf{Travaux anciens.} Oates, W. (1969), \emph{JPE} 77(6), 957-971 \lu{texte intégral} : sur
53 communes du New Jersey, environ deux tiers d'une hausse d'impôt foncier se retrouvent dans les
prix ; King (1977, \emph{JPE} 85(2)) ramène l'estimation à environ 40~\% en respécifiant le modèle.
Palmon, O. et B. Smith (1998), « New Evidence on Property Tax Capitalization », \emph{JPE} 106(5),
1099-1111 \lu{texte intégral} exploitent des districts de services de la banlieue de Houston
offrant les mêmes services à des taux d'impôt très différents, hérités de leur endettement : la
capitalisation est d'environ 62~\% avec un taux d'actualisation de 3~\%, complète avec 6,5~\%.
Le taux de capitalisation estimé dépend toujours du taux d'actualisation supposé, ce qui fragilise
ces résultats.

\textbf{Royaume-Uni.} Koster, H. et E. Pinchbeck (2022), « How Do Households Value the Future?
Evidence from Property Taxes », \emph{AEJ: Economic Policy} 14(1), 207-239 \lu{texte intégral,
version acceptée}. Ventes répétées de logements situés à moins de 1-2~km des frontières entre
collectivités votant des \emph{council taxes} différentes, environ vingt ans de transactions.
L'écart d'impôt est capitalisé comme si les acheteurs actualisaient au taux de 3,7~\% net de la
croissance attendue (« entre 3 et 4~\% »), proche des taux de marché : capitalisation quasi
complète ; la répercussion sur les loyers est d'environ 100~\%. Il s'agit d'un impôt sur le bien
entier, en montants forfaitaires par tranche de valeur.

\textbf{Danemark, première étude.} Høj, A., M. Jørgensen et P. Schou (2018), « Land Tax Changes and
Full Capitalisation », \emph{Fiscal Studies} 39(2), 365-380 \lu{version de travail 2017, texte
intégral}. La fusion des communes danoises en 2007 a obligé à harmoniser le taux de la taxe
municipale sur la valeur du sol : hausse dans 139 zones (+3,4~‰ en moyenne), baisse dans 105
($-2{,}6$~‰). Différence de différences à traitement continu sur 211~209 ventes de maisons
(2001-2008). Un point pour mille de taux effectif réduit le prix de 4,3~\% en 2007, ce qui implique
un taux d'actualisation de 2,3~\% : capitalisation complète. Les services publics ne sont
contrôlés que par un indicateur indirect.

\textbf{Danemark, seconde étude.} Nielsson, U., C. Wroblewski et A. Yding (2024), « The Incidence
and Efficiency of Land Value Taxation », document de travail, novembre 2024 \lu{texte intégral}.
Même réforme, deux sources de variation : le choc de taux entre communes voisines, et
l'exposition de chaque logement, mesurée par l'intensité foncière historique de sa parcelle (un
terrain nu est taxé sur toute sa valeur, un appartement très peu). Indice de ventes répétées,
effets fixes par logement et par nouvelle commune et année. Effet nul et précis sur les prix : la
borne inférieure de l'intervalle à 95~\% exclut une capitalisation complète pour tout taux
d'actualisation inférieur à 8~\%, alors que les estimations récentes se situent entre 2 et 5~\%.
Résultat identique sur les maisons individuelles. Pas d'effet sur la construction, la mise en
valeur des terrains vides, la mobilité ni l'accession ; les propriétaires âgés quittent toutefois
les zones où la taxe a augmenté. Les auteurs reproduisent le résultat de Høj et al. avec une
spécification proche de la leur et montrent qu'il disparaît quand on exclut les résidences de
vacances, plus nombreuses dans les zones de hausse et dont les prix ont divergé au même moment.
Interprétation : la charge est partagée avec les acheteurs futurs et les locataires, la taxe est
plus régressive que dans les modèles standard. Explication proposée : le contrôle des loyers
danois permet explicitement au bailleur de répercuter les hausses d'impôt. Elle n'explique pas
le résultat sur les maisons occupées par leur propriétaire, que les auteurs jugent spéculatif.
Étude non encore publiée.

\textbf{Répercussion sur les loyers.} Löffler, M. et S. Siegloch (2024), « Welfare Effects of
Property Taxation », ECONtribute DP 331 \lu{texte intégral ; versions IZA et CESifo de 2021}.
Environ 5~500 changements des multiplicateurs communaux de l'impôt foncier allemand (sur le
terrain et le bâti), surtout des hausses motivées par la consolidation budgétaire ; études
d'événement sur les loyers d'annonces. Un point de taux effectif accroît les loyers totaux de
3,4~\% après trois ans, soit une répercussion d'environ 83~\% (« complète » dans la version 2021) ;
peu d'effet sur salaires et profits. L'incidence est régressive : environ 1,3~\% de la consommation
pour le premier décile, 0,6~\% pour le dernier. Pas de publication en revue trouvée.

Baker, S. (2025), « Property Tax Pass-Through to Renters: A Quasi-Experimental Approach »,
Federal Reserve Bank of Philadelphia WP 25-41 \lu{introduction}. En Californie, la valeur
imposable est réinitialisée au prix d'achat à chaque vente ; à Berkeley, chaque dollar d'impôt
supplémentaire se retrouve pour 0,50 à 0,89 dollar dans le loyer des nouveaux locataires.

\section{Construction et densité}

Oates, W. et R. Schwab (1997), « The Impact of Urban Land Taxation: The Pittsburgh Experience »,
\emph{National Tax Journal} 50(1), 1-21 \lu{texte intégral}. En 1979-1980, Pittsburgh porte son
taux municipal sur le terrain à plus de cinq fois celui sur le bâti ; avec les taux du comté et des
écoles, uniformes, le rapport total n'est que d'environ deux. La valeur réelle des permis de
construire augmente de 70~\% entre 1960-1979 et 1980-1989, contre $-14$~\% en moyenne pour quinze
villes de la \emph{Rust Belt} ; Pittsburgh est la seule dont l'indicatrice post-réforme est
positive et significative. Les auteurs attribuent l'essentiel du boom à une pénurie de bureaux et
au programme Renaissance~II ; la taxe aurait joué un rôle d'appoint, en évitant d'autres hausses
d'impôt. Une seule ville traitée, sans contrefactuel construit.

Plassmann, F. et T.~N. Tideman (2000), « A Markov Chain Monte Carlo Analysis of the Effect of
Two-Rate Property Taxes on Construction », \emph{Journal of Urban Economics} 47(2), 216-247
\lu{résumé et synthèses}. Quinze communes de Pennsylvanie à deux taux comparées à 204 communes à
taux unique, 1972-1994 : davantage de permis de construire, de l'ordre de 3 à 4~\% pour le
résidentiel selon la synthèse de Dye et England (2010). Le résultat dépend de la spécification.

Banzhaf, H.~S. et N. Lavery (2010), « Can the Land Tax Help Curb Urban Sprawl? Evidence from
Growth Patterns in Pennsylvania », \emph{Journal of Urban Economics} 67(2), 169-179 \lu{version de
travail 2008, texte intégral}. Secteurs de recensement de Pennsylvanie, 1970-2000, premières
différences avec effets fixes et appariement. L'adoption du double taux accroît la croissance du
nombre de pièces de 5 à 6 points par décennie, surtout par davantage de logements plutôt que des
logements plus grands ; selon les modèles, +2 à +10 points sur la croissance de la densité de
logements. Les auteurs refusent de conclure à une réduction de l'étalement.

Dye, R. et R. England (2010), \emph{Assessing the Theory and Practice of Land Value Taxation},
Lincoln Institute \lu{texte intégral}. Synthèse : la LVT améliorerait l'équité et l'efficacité de
l'impôt foncier, mais les résultats statistiques sont « le plus souvent peu concluants ».
Recommandations : réévaluations fréquentes, mise en place progressive, crédit d'impôt pour les
ménages riches en terrain et pauvres en revenu. Pittsburgh a abandonné le double taux en 2001,
après une réévaluation chaotique.

Travaux récents : Hanson (2021) et Yang, Zhou et Hawley (2021), documents de travail du Lincoln
Institute \lu{résumés}, trouvent des effets positifs du double taux sur les implantations
d'entreprises et les valeurs agrégées ; ils ne sont pas publiés en revue à notre connaissance.

\section{Droits de mutation}

Hilber, C. et T. Lyytikäinen (2017), « Transfer taxes and household mobility: Distortion on the
housing or labor market? », \emph{Journal of Urban Economics} 101, 57-73 \lu{texte intégral}.
Discontinuité du \emph{stamp duty} britannique à 250~000~£ (de 1 à 3~\%) : 2 points de droit en
plus réduisent la mobilité annuelle d'environ 37~\%, surtout les déménagements courts liés au
logement, sans effet sur la mobilité professionnelle. Perte sèche de 40 à 80~\% de la recette
supplémentaire.

Best, M. et H. Kleven (2018), « Housing Market Responses to Transaction Taxes: Evidence From
Notches and Stimulus in the U.K. », \emph{Review of Economic Studies} 85(1), 157-193 \lu{résumé}.
Les seuils du \emph{stamp duty} provoquent des distorsions fortes des prix et des volumes ; la
suspension temporaire d'un droit de 1~\% en 2008-2009 accroît l'activité d'environ 20~\% à court
terme.

Dachis, B., G. Duranton et M. Turner (2012), « The effects of land transfer taxes on real estate
markets: evidence from a natural experiment in Toronto », \emph{Journal of Economic Geography}
12(2), 327-354 \lu{résumé et version de travail}. Un droit municipal de 1,1~\% réduit les ventes de
15~\% ; un dollar de perte sèche pour huit dollars de recette, comparé à un impôt foncier de même
rendement.

Fritzsche, C. et L. Vandrei (2019), \emph{Regional Science and Urban Economics} 74, 131-143
\lu{résumé} : un point de droit supplémentaire réduit les transactions de maisons d'environ 7~\% en
Allemagne, avec des ventes avancées avant chaque hausse. Dolls, M., C. Fuest, C. Krolage et
F. Neumeier (2025), « Who bears the burden of real estate transfer taxes? », \emph{Journal of
Urban Economics} 145, 103717 \lu{résumé} : sur environ 17 millions de biens (2005-2019), un point de droit réduit les
prix de 3~\% après un an, plus que la charge d'une seule transaction, ce que les auteurs
attribuent à la capitalisation des droits des reventes futures.

Bontemps, C., F. Cherbonnier et T. Magnac (2023), « Reducing transaction taxes on housing in highly
regulated economies », TSE WP 23-1486 \lu{texte intégral, passages principaux}. Modèle à
générations et agents hétérogènes, avec contraintes de crédit (apport de 20~\%, taux d'effort de
30~\%), estimé sur l'Enquête Logement 2006. Ramener les coûts de transaction de 8~\% à l'achat et
6~\% à la vente à 3~\% et 3~\%, à budget équilibré par un impôt proportionnel sur la propriété :
transactions +21~\%, accès à la propriété élargi aux jeunes aisés et réduit pour les jeunes
modestes, prix +2,4~\%, loyers $-1{,}6$~\%, bien-être +0,2~\%. Le gain est faible parce que les
contraintes de crédit, plus que les droits, bornent les choix.

\section{Faits et institutions : France et comparaisons}

\textbf{Le prix du sol.} Knoll, K., M. Schularick et T. Steger (2017), « No Price Like Home: Global
House Prices, 1870-2012 », \emph{AER} 107(2), 331-353 \lu{résumé publié, version de travail
intégrale}. Prix réels des logements dans quatorze pays : quasi constants jusqu'au milieu du
XX\textsuperscript{e} siècle, puis multipliés ; la hausse du prix du sol explique « environ
80~\% » de la hausse depuis 1950. La version de travail donne 93~\% pour la France, chiffre non
vérifié dans la version publiée. Décomposition comptable, sensible à la part supposée du sol.

\textbf{Régressivité de la taxe foncière.} André, M. et O. Meslin (2023), « Les ménages au
patrimoine immobilier le plus élevé paient relativement moins de taxe foncière que les ménages les
moins dotés », \emph{Insee Analyses} n°~91 \lu{texte intégral et données}. Appariement des
fichiers fiscaux de 2017 avec une valeur de marché de chaque logement estimée sur les ventes
2015-2019. La taxe représente 0,34~\% du patrimoine immobilier brut en moyenne, plus de 0,5~\%
pour les plus petits patrimoines, 0,22~\% pour le centile supérieur, 0,18~\% pour le millième
supérieur ; environ 0,1~\% de la valeur à Paris, souvent plus de 0,4~\% en Seine-Saint-Denis.
Rapportée au revenu, elle est progressive. Cause principale : l'écart entre valeurs locatives de
1970 et valeurs de marché.

\textbf{Rapports publics.} Conseil des prélèvements obligatoires (2023), \emph{Pour une fiscalité
du logement plus cohérente} \lu{texte intégral} : la fiscalité du logement représente 92~Md€ en
2022 ; la base de la taxe foncière est « largement décorrélée » des loyers et des prix ; le CPO
recommande une base liée aux loyers ou aux prix de marché, sans exclure la valeur vénale, puis
une baisse des droits de mutation compensée par l'impôt sur la détention. Chapelle, G., B. Fabre
et C. Lallemand (2020), rapport IPP n°~29 \lu{texte intégral} : la révision des valeurs
locatives d'habitation sur les loyers augmenterait celles des logements anciens et des ménages
aisés, baisserait celles des logements des années 1960-1970 et des ménages modestes. Trannoy, A.
et É. Wasmer (2013), « Comment modérer les prix de l'immobilier ? », \emph{Note du CAE} n°~2
\lu{texte intégral} : supprimer progressivement les droits de mutation et asseoir la taxe foncière
sur la valeur vénale nette des dettes ; taxer la plus-value latente quand un terrain devient
constructible. Leur livre \emph{Le grand retour de la terre dans les patrimoines} (Odile Jacob,
2022) propose un impôt plat sur la valeur des terrains \lu{recensions seulement ; taux proposé non
vérifié}.

\textbf{Synthèses internationales.} Mirrlees, J. et al. (2011), \emph{Tax by Design}, ch.~16
\lu{texte intégral} : l'argument en faveur d'une LVT est « simple et presque indéniable » ; la
valeur n'a pas besoin d'être exacte pour que la taxe soit efficace. Proposition : remplacer les
\emph{business rates} par une LVT sur les terrains professionnels et agricoles (environ 4~\% de la
valeur, introduite à partir de 0,5~\%), et la \emph{council tax} et les droits de mutation par un
impôt sur la valeur des services de logement, sans LVT résidentielle séparée. OCDE (2022),
\emph{Housing Taxation in OECD Countries} \lu{texte intégral} : seuls la Nouvelle-Galles du Sud, le
Danemark et l'Estonie ont une taxe sur le seul terrain ; recommandation de renforcer l'impôt
récurrent sur des valeurs actualisées et de réduire les droits de mutation.

\textbf{Allemagne.} Après la censure par la Cour constitutionnelle (10 avril 2018) d'un impôt
foncier assis sur des valeurs de 1964 et 1935, la réforme de 2019, appliquée en 2025, a laissé les
Länder choisir. Le Bade-Wurtemberg taxe le seul terrain : surface × valeur de zone du sol, taux de
base de 1,30~‰ réduit de 30~\% pour l'usage résidentiel, preuve d'une valeur inférieure admise
au-delà de 30~\% d'écart. Le modèle fédéral a été validé par la Cour fédérale des finances en
décembre 2025 (II~R~25/24 et autres), celui du Bade-Wurtemberg le 20~mai 2026 (II~R~26/24 et
II~R~27/24) \lu{communiqués de la Cour}. Aucune étude académique des effets n'est disponible ; les
chiffres circulant sur le transfert de charge vers le logement viennent d'associations.

\section*{Ce qui manque}

Aucune étude causale n'existe sur la France, faute de réforme exploitable. La question décisive,
savoir si une taxe sur le seul sol est capitalisée ou répercutée, repose sur une seule réforme
étudiée deux fois avec des conclusions opposées. La réforme du Bade-Wurtemberg, qui a modifié en
2025 la charge relative de millions de parcelles selon la part du terrain dans leur valeur, offre
la meilleure occasion de trancher dans les années qui viennent.

\end{document}
